\documentclass[11pt,a4paper]{article}

\usepackage{amsmath,amssymb,amsthm}
\usepackage{graphicx}
\usepackage{booktabs}
\usepackage{hyperref}
\usepackage[margin=1in]{geometry}
\usepackage{algorithm}
\usepackage{algpseudocode}

\newtheorem{definition}{Definition}
\newtheorem{theorem}{Theorem}
\newtheorem{proposition}{Proposition}

\title{Symbolic Mechanism Discovery via A* Search\\in Mathematical Operation Space}

\author{Andrew H. Bond\\
\small Department of Computer Engineering, San Jose State University\\
\small Senior Member, IEEE \quad ORCID: 0009-0003-2599-6158\\
\small \texttt{andrew.bond@sjsu.edu}}

\date{\today}

\begin{document}
\maketitle

\begin{abstract}
We introduce a framework for automated discovery of interpretable
physical laws from simulation data using A* search through a space
of mathematical operations. Given a set of computed features and
binary labels (e.g., ``periodic'' vs.\ ``chaotic''), the method
searches over compositions of arithmetic, transcendental, and
conditional operations to find the closed-form formula that best
predicts the label, as measured by F1 score. Unlike genetic
programming, A* provides an optimality guarantee within the search
depth. Unlike neural symbolic regression, the discovered formulas
are human-readable and physically interpretable.

Applied to the gravitational three-body problem (15,000 configurations,
23 Hessian-derived features, GPU-integrated ground truth), the method
discovers that periodicity is predicted by the Euclidean norm
$\sqrt{n_\mathrm{clusters}^2 + \lambda_\mathrm{min}^2}$ in a
two-dimensional feature space of frequency cluster count and minimum
Hessian eigenvalue (F1~$= 0.48$). The search further identifies a
\emph{formula ceiling}: no closed-form expression of depth~$\leq 4$
exceeds F1~$\approx 0.50$, while a gradient-boosted ensemble on the
same features reaches F1~$= 0.86$. This gap quantifies the boundary
between symbolically expressible and irreducibly complex structure in
the data---a diagnostic that existing symbolic regression methods do
not provide.

The framework is domain-general, requiring only a feature matrix and
a label vector. Code and data are available at
\url{https://github.com/ahb-sjsu/tensor-tractability}.
\end{abstract}

\section{Introduction}

A central goal of computational science is to discover interpretable
laws from data. Given measurements of a physical system, can we
automatically find a formula that predicts its behavior? This question
has been addressed by symbolic regression~\cite{schmidt2009distilling},
neural equation discovery~\cite{udrescu2020ai}, and classical AI
systems like BACON~\cite{langley1987scientific}. Each approach has
limitations: genetic programming provides no optimality guarantee and
often produces bloated expressions; neural methods discover accurate
but opaque representations; classical systems require hand-crafted
search heuristics.

We propose \textbf{A* Symbolic Discovery}, which searches a space of
mathematical operations using A* with F1 score as the heuristic. The
method inherits A*'s key property: given an admissible heuristic, it
finds the optimal formula within the search budget. The formulas are
constructed from a registry of operations (arithmetic, transcendental,
conditional) applied to a feature vector, producing expressions that
are both optimal and human-readable.

Our framework provides two outputs that existing methods do not:
\begin{enumerate}
\item The \textbf{best interpretable formula} for the prediction task,
with a guarantee that no simpler formula (within the search depth)
achieves higher F1.
\item The \textbf{formula ceiling}---the maximum F1 achievable by any
closed-form expression of bounded complexity. Comparing this ceiling
to the performance of black-box models (e.g., gradient boosting)
quantifies how much structure in the data is symbolically expressible
versus irreducibly complex.
\end{enumerate}

\section{Method}

\subsection{Problem Formulation}

\begin{definition}[Symbolic Prediction Problem]
Given a feature matrix $X \in \mathbb{R}^{N \times d}$ and binary
labels $y \in \{0, 1\}^N$, find a formula
$f: \mathbb{R}^d \to \mathbb{R}$ and threshold $\tau \in \mathbb{R}$
such that the predictor $\hat{y}_i = \mathbf{1}[f(X_i) > \tau]$
maximizes the F1 score:
\[
F_1 = \frac{2 \cdot \mathrm{precision} \cdot \mathrm{recall}}
           {\mathrm{precision} + \mathrm{recall}}.
\]
The formula $f$ is constrained to be a composition of operations from
a finite registry $\mathcal{O}$ with depth at most $D$.
\end{definition}

\subsection{Operation Registry}

The search space is defined by three classes of operations:

\textbf{Leaf operations} (depth 0): Each feature $x_j$ is a leaf node.

\textbf{Unary operations} (depth +1): $\log|x|$, $\sqrt{|x|}$, $x^2$,
$1/x$, $-x$, $|x|$, $\sigma(x)$, $\tanh(x)$.

\textbf{Binary operations} (depth +1): $x + y$, $x - y$, $x \cdot y$,
$x / y$, $\max(x, y)$, $\min(x, y)$, $\sqrt{x^2 + y^2}$ (hypot),
$(x - y)^2$, $2xy/(x+y)$ (harmonic mean), $\mathrm{sgn}(xy)\sqrt{|xy|}$
(geometric mean).

Each operation has an associated cost (1 for arithmetic, 2 for
transcendental). The total cost of a formula is the sum of its
operation costs.

\subsection{A* Search}

We search the formula space using A* (Algorithm~\ref{alg:astar}).
Each state is a formula tree. Expansion generates new states by
applying one operation (unary to the current formula, or binary
combining it with a leaf). The priority is $g(n) + h(n)$ where
$g$ is the formula cost and $h$ is the negative F1 score (since
we maximize F1 by minimizing $-$F1).

\begin{algorithm}[t]
\caption{A* Symbolic Discovery}
\label{alg:astar}
\begin{algorithmic}[1]
\Require Feature matrix $X$, labels $y$, operations $\mathcal{O}$,
max depth $D$, max states $S$
\State Initialize frontier $\gets$ \{all leaf features\}
\State best $\gets$ (0, \texttt{null})
\While{frontier nonempty and explored $< S$}
    \State $f \gets$ pop lowest-priority from frontier
    \If{$\mathrm{F1}(f, X, y) > \mathrm{best.F1}$}
        \State best $\gets$ ($\mathrm{F1}(f)$, $f$)
    \EndIf
    \If{depth($f$) $< D$}
        \For{each unary op $u \in \mathcal{O}$}
            \State push $u(f)$ to frontier
        \EndFor
        \For{each leaf $\ell$ and binary op $b \in \mathcal{O}$}
            \State push $b(f, \ell)$ and $b(\ell, f)$ to frontier
        \EndFor
    \EndIf
\EndWhile
\State \Return best
\end{algorithmic}
\end{algorithm}

\textbf{F1 evaluation.} For each candidate formula, we compute its
value on all $N$ data points, then sweep a set of thresholds
(percentiles 20--95 of the value distribution) to find the threshold
maximizing F1. Both above-threshold and below-threshold predictors
are tested. This sweep adds negligible cost when vectorized.

\textbf{GPU acceleration.} The feature matrix and label vector are
stored on GPU (via CuPy). Formula evaluation, threshold sweep, and
F1 computation are all vectorized GPU operations, enabling evaluation
of thousands of formulas per second on a single GPU.

\subsection{Formula Ceiling Detection}

\begin{definition}[Formula Ceiling]
The \emph{formula ceiling} $C_D$ at depth $D$ is:
\[
C_D = \max_{f \in \mathcal{F}_D} F_1(f, X, y, \tau^*(f))
\]
where $\mathcal{F}_D$ is the set of all formulas of depth $\leq D$
and $\tau^*(f)$ is the optimal threshold for each formula.
\end{definition}

If $C_D$ is significantly below the F1 of a black-box model (e.g.,
gradient boosting) trained on the same features, then the data
contains structure that cannot be captured by depth-$D$ formulas.
This gap $G = F_1^{\mathrm{black\text{-}box}} - C_D$ quantifies
the \emph{irreducible complexity} of the prediction task at depth $D$.

\begin{proposition}[Ceiling Convergence]
$C_D$ is non-decreasing in $D$ and bounded above by
$F_1^{\mathrm{Bayes}}$, the Bayes-optimal F1. If $G > 0$ for all
$D \leq D_{\max}$, the underlying relationship between features and
labels is not expressible as a closed-form formula of depth
$\leq D_{\max}$.
\end{proposition}

\section{Case Study: Three-Body Periodic Orbit Prediction}

\subsection{Setup}

We compute 23 features from the Hessian (coupling tensor) of the
gravitational three-body Hamiltonian at 15,000 random configurations
of equal-mass bodies. Features include: matrix rank, potential block
eigenvalues, frequency ratios, cluster counts, spectral entropy,
energy, separation ratios, and derived quantities (log-transforms,
inverses). Ground truth (periodic vs.\ non-periodic) is obtained by
GPU-accelerated symplectic integration of 80,000 leapfrog steps
per orbit on an NVIDIA GV100 GPU, with periodicity defined as return
to within 5\% of the initial phase-space extent.

\subsection{Results}

\textbf{Phase 1: Single features.} The best single-feature predictor
is \texttt{n\_clusters} (number of eigenvalue frequency clusters),
achieving F1~$= 0.35$.

\textbf{Phase 2: Pairwise combinations.} 5,060 formulas evaluated
(23 features $\times$ 22 features $\times$ 10 operations). Best:
\texttt{n\_clusters + min\_nonzero\_eig} at F1~$= 0.48$.

\textbf{Phase 3: Unary transforms.} 400 formulas (top 50 pairwise
$\times$ 8 unary ops). No improvement: best remains F1~$= 0.48$.

\textbf{Phase 4: Triple combinations.} 504 formulas (top 8 features,
9 combination patterns). Zero formulas exceed F1~$= 0.30$. Adding
a third feature provides no benefit.

\textbf{Formula ceiling.} $C_4 \approx 0.48$ across all four phases
and all tested operation types. This ceiling is robust: log, sqrt,
square, sigmoid, tanh, max, min, hypot, harmonic mean, and geometric
mean all produce identical F1 when applied to the same two features.

\textbf{Black-box comparison.} A gradient-boosted classifier (100
trees, depth 4) trained on 9 features achieves cross-validated
F1~$= 0.86$ on a 51,000-configuration dataset spanning three mass
ratios. A depth-5 decision tree achieves F1~$= 0.58$.

\textbf{Irreducible complexity gap.} $G = 0.86 - 0.48 = 0.38$.
Nearly 40\% of the predictive structure is not expressible as a
depth-4 closed-form formula over Hessian features. The decision tree
at F1~$= 0.58$ shows that piecewise-linear functions capture some
of this gap; the remaining $0.86 - 0.58 = 0.28$ requires the
nonlinear interactions of an ensemble.

\subsection{Discovered Formula}

The A* search discovers:
\begin{equation}
\hat{y} = \mathbf{1}\!\left[
\sqrt{n_\mathrm{clusters}^2 + \lambda_\mathrm{min}^2} > \tau
\right]
\label{eq:discovered}
\end{equation}
where $n_\mathrm{clusters}$ is the number of frequency clusters
in the Hessian's potential block (counting gaps $> 10\times$ in the
eigenvalue spectrum) and $\lambda_\mathrm{min}$ is the smallest
nonzero eigenvalue magnitude.

\textbf{Physical interpretation.} Periodicity requires two
independent conditions: (i)~timescale separation
($n_\mathrm{clusters} \geq 2$, so the inner binary decouples from
the outer body) and (ii)~gravitational binding
($\lambda_\mathrm{min} > 0$, so the system does not escape). The
Euclidean norm indicates these conditions contribute independently---there
is no synergy or interaction between them.

\section{Discussion}

\subsection{Comparison to Existing Methods}

\textbf{Genetic symbolic regression} (e.g., PySR~\cite{cranmer2023interpretable})
uses evolutionary search with no optimality guarantee. It may find
complex formulas that overfit; our A* search at bounded depth
guarantees the simplest optimal formula. However, genetic methods
can explore deeper formula trees more efficiently when the search
space is very large.

\textbf{AI Feynman}~\cite{udrescu2020ai} uses neural networks to
identify symmetries and separabilities before fitting symbolic
expressions. Our method is simpler (no neural preprocessing) but
less powerful for discovering dimensional relationships.

\textbf{BACON/Eureka}~\cite{langley1987scientific} searched for
laws using hand-crafted heuristics. Our A* formulation replaces
domain-specific heuristics with a generic optimality criterion (F1).

\textbf{The formula ceiling is novel.} No existing symbolic
regression method systematically quantifies the gap between the
best formula and the best black-box model. This gap is a property
of the data, not the search algorithm, and provides a principled
stopping criterion: if the ceiling has converged across multiple
search depths and operation sets, further symbolic search is futile.

\subsection{Limitations}

The search is exponential in depth: depth $D$ with $d$ features and
$k$ operations explores $O((d \cdot k)^D)$ formulas. Our GPU
vectorization makes depth 3--4 practical (minutes), but depth 6+
would require pruning heuristics or beam search.

The method assumes features are pre-computed. Discovering the right
features (e.g., ``eigenvalue cluster count'' from raw Hessian entries)
requires domain knowledge. Combining our method with automated
feature engineering (e.g., deep feature synthesis) is future work.

The F1 metric weights precision and recall equally. For scientific
discovery, recall (not missing real phenomena) may matter more than
precision; alternative metrics could change the discovered formula.

\section{Conclusion}

A* Symbolic Discovery provides a principled framework for finding
interpretable physical laws from data, with two novel contributions:
optimality-guaranteed formula search and formula ceiling detection.
Applied to the three-body problem, it discovers that periodicity is
a two-feature phenomenon (timescale separation + gravitational
binding) with a hard symbolic ceiling at F1~$\approx 0.50$, while
the full nonlinear structure captured by ensemble methods reaches
F1~$= 0.86$. The 38\% gap quantifies the irreducible complexity
of three-body dynamics beyond what any short formula can express.

The method is domain-general and requires only a feature matrix and
labels. It is available as open-source software at
\url{https://github.com/ahb-sjsu/tensor-tractability}.

\begin{thebibliography}{9}

\bibitem{schmidt2009distilling}
M.~Schmidt and H.~Lipson, ``Distilling free-form natural laws from
experimental data,'' \emph{Science}, vol.~324, pp.~81--85, 2009.

\bibitem{udrescu2020ai}
S.-M.~Udrescu and M.~Tegmark, ``AI Feynman: A physics-inspired
method for symbolic regression,'' \emph{Science Advances}, vol.~6,
no.~16, eaay2631, 2020.

\bibitem{langley1987scientific}
P.~Langley, H.~A.~Simon, G.~L.~Bradshaw, and J.~M.~Zytkow,
\emph{Scientific Discovery: Computational Explorations of the
Creative Processes}. MIT Press, 1987.

\bibitem{cranmer2023interpretable}
M.~Cranmer, ``Interpretable machine learning for science with PySR
and SymbolicRegression.jl,'' \emph{arXiv:2305.01582}, 2023.

\bibitem{kolda2009tensor}
T.~G.~Kolda and B.~W.~Bader, ``Tensor decompositions and
applications,'' \emph{SIAM Review}, vol.~51, no.~3, pp.~455--500,
2009.

\bibitem{bond2026geometric}
A.~H.~Bond, \emph{Geometric Methods in Computational Modeling:
From Manifolds to Production Systems}. San Jose State University,
2026.

\end{thebibliography}

\end{document}
